\chapter{Технологическая часть}

В данном разделе будут приведены средства реализации, листинг кода и функциональные тесты.

\section{Средства реализации}

В качестве языка программирования для реализации данной лабораторной работы был выбран язык программирования \texttt{C}~\cite{c-lang}

В этой лабораторной работе необходимо измерить время выполнения процессора для каждой реализации алгоритма и построить графики, показывающие зависимость времени выполнения от длины строк. Все необходимые инструменты для выполнения этой задачи доступны в выбранном языке программирования.

Время работы замерено с помощью функции \texttt{clock()} из библиотеки \texttt{time}~\cite{c-lang-time}.

\section{Сведения о модулях программы}

Данная программа разбита на следующие модули:

\begin{enumerate}
	\item \texttt{main.c}~--- главный файл программы содержащий точку входа в программу, где располагается вывод меню.
	\item \texttt{alloc\_matrix.c}~--- файл, содержащий реализацию динамической матрицы.
	\item \texttt{algorithms.c}~--- файл с реализацией исследуемых алгоритмов.
	\item \texttt{measurement.c}~--- файл с функциями, замеряющими время выполнения алгоритмов поиска расстояния Левенштейна и Дамерау --- Левенштейна.  
\end{enumerate}

\section{Реализация алгоритмов}

В листингах~\ref{lst:lr.c}~--~\ref{lst:dlnr.c} в приложении~A приведены реализации алгоритмов поиска расстояний Левенштейна и Дамерау~---~Левенштейна.

\section{Функциональные тесты}

В таблице~\ref{tbl:func_tests} представлены функциональные тесты для реализованных алгоритмов. Тестирование проведено по методу чёрного ящика. Все тесты пройдены успешно.
\begin{table}[ht]
	\small
	\begin{center}
		\begin{threeparttable}
		\caption{Функциональные тесты}
		\label{tbl:func_tests}
		\begin{tabular}{|c|c|c|c|c|c|}
			\hline
			\multicolumn{3}{|c|}{\bfseries Входные данные} 
			& \multicolumn{2}{c|}{\bfseries Ожидаемый результат} \\ 
			\hline
			№ & Строка 1 & Строка 2 & Левенштейн & Дамерау~---~Левенштейн \\
			\hline
			1 & ab & ab & 0 & 0 \\
			\hline
			2 & aa & df & 2 & 2 \\
			\hline
			3 & asdfg & sof & 3 & 3 \\
			\hline
			4 & string & slakng & 3 & 3 \\
			\hline
			5 & dog & dooog & 2 & 2 \\
			\hline
			6 & dog & doggy & 2 & 2 \\
			\hline
			7 & dog & god & 2 & 2 \\
			\hline
			8 & "пустая строка" & "пустая строка" & 0 & 0 \\
			\hline
			9 & lack & nolack & 2 & 2 \\
			\hline
			10 & kronen & lampad & 6 & 6 \\
			\hline
		\end{tabular}	
		\end{threeparttable}
	\end{center}
\end{table}

\section*{Вывод}
В данном были разработаны и протестированы реализации алгоритмов нахождения расстояния Левенштейна и Дамерау~---~Левенштейна.